\section{Coalescing、tiling 与矩阵性能异常}

本节解释 DRAM burst、memory coalescing、row-major matmul、tiling、alignment、wave quantization 和 matrix mystery。


\begin{knowledgebox}{first-use glossary：coalescing、tiling、wave quantization}
Memory coalescing 指同一个 warp 的 memory accesses 落在连续 burst 中，从而合并成高效访问。Tiling 把大矩阵切成小块放入 shared memory 重复使用。Wave quantization 指 thread blocks 数量不能完美填满 SM waves 时产生的周期性利用率波动。
\end{knowledgebox}

\begin{figure}[H]
\centering
\includegraphics[width=0.88\textwidth]{slide-036.jpg}
\caption{Slide 36: Throw away the activations, re-compute them!}
\figsource{CS336 Spring 2026 Lecture 5 官方幻灯片。}
\end{figure}

\noindent\textbf{展开说明：}Throwing away computation can actually be optimal, w/ 5/8 th the memory accesses!

\begin{knowledgebox}{读图：Slide 36 应该怎么看}
这页是硬件机制或性能证据页。先看数据从哪里来、在哪里复用、是否写回 HBM；再判断瓶颈是 compute、memory bandwidth、thread scheduling 还是 alignment。GPU 优化不是让公式更漂亮，而是让数据在正确层级停留更久、让 tensor cores 少等数据。
\end{knowledgebox}


\begin{figure}[H]
\centering
\includegraphics[width=0.88\textwidth]{slide-037.jpg}
\caption{Slide 37: Trick (?) 4: Memory coalescing and DRAM}
\figsource{CS336 Spring 2026 Lecture 5 官方幻灯片。}
\end{figure}

\noindent\textbf{展开说明：}DRAM (global memory) is read in ‘burst mode’ – each read gives you many bytes!

\begin{knowledgebox}{读图：Slide 37 应该怎么看}
这页是硬件机制或性能证据页。先看数据从哪里来、在哪里复用、是否写回 HBM；再判断瓶颈是 compute、memory bandwidth、thread scheduling 还是 alignment。GPU 优化不是让公式更漂亮，而是让数据在正确层级停留更久、让 tensor cores 少等数据。
\end{knowledgebox}


\begin{figure}[H]
\centering
\includegraphics[width=0.88\textwidth]{slide-038.jpg}
\caption{Slide 38: Memory coalescing}
\figsource{CS336 Spring 2026 Lecture 5 官方幻灯片。}
\end{figure}

\noindent\textbf{展开说明：}Memory accesses are coalesced if all the threads (in a warp) fall within the same burst

\begin{knowledgebox}{读图：Slide 38 应该怎么看}
这页是硬件机制或性能证据页。先看数据从哪里来、在哪里复用、是否写回 HBM；再判断瓶颈是 compute、memory bandwidth、thread scheduling 还是 alignment。GPU 优化不是让公式更漂亮，而是让数据在正确层级停留更久、让 tensor cores 少等数据。
\end{knowledgebox}


\begin{figure}[H]
\centering
\includegraphics[width=0.88\textwidth]{slide-039.jpg}
\caption{Slide 39: Coalescing for matrix multiplication}
\figsource{CS336 Spring 2026 Lecture 5 官方幻灯片。}
\end{figure}

\noindent\textbf{展开说明：}For row-major matrices – threads that move along rows are not coalesced

\begin{knowledgebox}{读图：Slide 39 应该怎么看}
这页是硬件机制或性能证据页。先看数据从哪里来、在哪里复用、是否写回 HBM；再判断瓶颈是 compute、memory bandwidth、thread scheduling 还是 alignment。GPU 优化不是让公式更漂亮，而是让数据在正确层级停留更久、让 tensor cores 少等数据。
\end{knowledgebox}


\begin{figure}[H]
\centering
\includegraphics[width=0.88\textwidth]{slide-040.jpg}
\caption{Slide 40: Trick 5 (the big one): tiling}
\figsource{CS336 Spring 2026 Lecture 5 官方幻灯片。}
\end{figure}

\noindent\textbf{展开说明：}Tiling is the idea of grouping and ordering threads to minimize global memory access.

\begin{knowledgebox}{读图：Slide 40 应该怎么看}
这页是硬件机制或性能证据页。先看数据从哪里来、在哪里复用、是否写回 HBM；再判断瓶颈是 compute、memory bandwidth、thread scheduling 还是 alignment。GPU 优化不是让公式更漂亮，而是让数据在正确层级停留更久、让 tensor cores 少等数据。
\end{knowledgebox}


\begin{figure}[H]
\centering
\includegraphics[width=0.88\textwidth]{slide-041.jpg}
\caption{Slide 41: Tiling – store and reuse information in shared memory}
\figsource{CS336 Spring 2026 Lecture 5 官方幻灯片。}
\end{figure}

\noindent\textbf{展开说明：}Cut up the matrix into smaller ‘tiles’, and load this into shared memory

\begin{knowledgebox}{读图：Slide 41 应该怎么看}
这页是硬件机制或性能证据页。先看数据从哪里来、在哪里复用、是否写回 HBM；再判断瓶颈是 compute、memory bandwidth、thread scheduling 还是 alignment。GPU 优化不是让公式更漂亮，而是让数据在正确层级停留更久、让 tensor cores 少等数据。
\end{knowledgebox}


\begin{figure}[H]
\centering
\includegraphics[width=0.88\textwidth]{slide-042.jpg}
\caption{Slide 42: Tiling math}
\figsource{CS336 Spring 2026 Lecture 5 官方幻灯片。}
\end{figure}

\noindent\textbf{展开说明：}Non-tiled matrix multiply: each input is read $N$ times from global memory

\begin{knowledgebox}{读图：Slide 42 应该怎么看}
这页是硬件机制或性能证据页。先看数据从哪里来、在哪里复用、是否写回 HBM；再判断瓶颈是 compute、memory bandwidth、thread scheduling 还是 alignment。GPU 优化不是让公式更漂亮，而是让数据在正确层级停留更久、让 tensor cores 少等数据。
\end{knowledgebox}


\begin{figure}[H]
\centering
\includegraphics[width=0.88\textwidth]{slide-043.jpg}
\caption{Slide 43: Complexities with tiling}
\figsource{CS336 Spring 2026 Lecture 5 官方幻灯片。}
\end{figure}

\noindent\textbf{展开说明：}Tile sizes may not divide the matrix size and lead to low utilization

\begin{knowledgebox}{读图：Slide 43 应该怎么看}
这页是硬件机制或性能证据页。先看数据从哪里来、在哪里复用、是否写回 HBM；再判断瓶颈是 compute、memory bandwidth、thread scheduling 还是 alignment。GPU 优化不是让公式更漂亮，而是让数据在正确层级停留更久、让 tensor cores 少等数据。
\end{knowledgebox}


\begin{figure}[H]
\centering
\includegraphics[width=0.88\textwidth]{slide-044.jpg}
\caption{Slide 44: Complexities with tiling 2 – memory alignment}
\figsource{CS336 Spring 2026 Lecture 5 官方幻灯片。}
\end{figure}

\noindent\textbf{展开说明：}Memory comes in bursts

\begin{knowledgebox}{读图：Slide 44 应该怎么看}
这页是硬件机制或性能证据页。先看数据从哪里来、在哪里复用、是否写回 HBM；再判断瓶颈是 compute、memory bandwidth、thread scheduling 还是 alignment。GPU 优化不是让公式更漂亮，而是让数据在正确层级停留更久、让 tensor cores 少等数据。
\end{knowledgebox}


\begin{figure}[H]
\centering
\includegraphics[width=0.88\textwidth]{slide-045.jpg}
\caption{Slide 45: Putting it together: understanding a matrix mystery}
\figsource{CS336 Spring 2026 Lecture 5 官方幻灯片。}
\end{figure}

\noindent\textbf{展开说明：}Why is it faster to have bigger matrices?

\begin{knowledgebox}{读图：Slide 45 应该怎么看}
这页是硬件机制或性能证据页。先看数据从哪里来、在哪里复用、是否写回 HBM；再判断瓶颈是 compute、memory bandwidth、thread scheduling 还是 alignment。GPU 优化不是让公式更漂亮，而是让数据在正确层级停留更久、让 tensor cores 少等数据。
\end{knowledgebox}


\begin{figure}[H]
\centering
\includegraphics[width=0.88\textwidth]{slide-046.jpg}
\caption{Slide 46: Matrix mystery}
\figsource{CS336 Spring 2026 Lecture 5 官方幻灯片。}
\end{figure}

\noindent\textbf{展开说明：}We understand some of this (compute intensity, tiling). Let’s take a closer look..

\begin{knowledgebox}{读图：Slide 46 应该怎么看}
这页是硬件机制或性能证据页。先看数据从哪里来、在哪里复用、是否写回 HBM；再判断瓶颈是 compute、memory bandwidth、thread scheduling 还是 alignment。GPU 优化不是让公式更漂亮，而是让数据在正确层级停留更久、让 tensor cores 少等数据。
\end{knowledgebox}


\begin{figure}[H]
\centering
\includegraphics[width=0.88\textwidth]{slide-047.jpg}
\caption{Slide 47: Part 1: tiling}
\figsource{CS336 Spring 2026 Lecture 5 官方幻灯片。}
\end{figure}

\noindent\textbf{展开说明：}Tiling has a major impact through alignment.

\begin{knowledgebox}{读图：Slide 47 应该怎么看}
这页是硬件机制或性能证据页。先看数据从哪里来、在哪里复用、是否写回 HBM；再判断瓶颈是 compute、memory bandwidth、thread scheduling 还是 alignment。GPU 优化不是让公式更漂亮，而是让数据在正确层级停留更久、让 tensor cores 少等数据。
\end{knowledgebox}


\begin{figure}[H]
\centering
\includegraphics[width=0.88\textwidth]{slide-048.jpg}
\caption{Slide 48: Part 2: wave quantization}
\figsource{CS336 Spring 2026 Lecture 5 官方幻灯片。}
\end{figure}

\noindent\textbf{展开说明：}What’s with the periodic behavior?

\subsection{本章小结}
本节的共同问题是如何减少昂贵的数据移动并提高硬件利用率。GPU 的速度来自海量并行和专用矩阵硬件，但只有当 memory access、shape、precision 和 kernel 组织匹配时，这些速度才会真正出现。

